% =====================================================================
%  Section: Experiments  (IJCAI 2027 revision: IJCAI-style with figure
%  placeholders for qualitative results)
% =====================================================================
\section{Experiments}
\label{sec:experiment}

We design experiments to answer six questions:
\textbf{Q1}~Does Canvas3D-v2 outperform state-of-the-art image-to-video,
3D reconstruction, and painting-animation methods on a curated painting
benchmark?
\textbf{Q2}~How much does each of the seven modules contribute to the
final quality, and which design choices within each module are
load-bearing?
\textbf{Q3}~Do human viewers prefer Canvas3D-v2 outputs on artistic
authenticity, motion naturalness, and compositional-intent alignment?
\textbf{Q4}~What is the per-module compute cost, and where does the
runtime budget go?
\textbf{Q5}~How well does ViewCtrl support user-specified viewpoint
trajectories, and does the generated video maintain 3D consistency
across novel viewpoints?
\textbf{Q6}~How accurately does SpatialFlow preserve user-edited spatial
relations across generated frames, and how much does the relation
control signal contribute?

\subsection{Experimental Setup}

\paragraph{Dataset.}
We curate the \textbf{PaintingAnim-50} benchmark of 50 high-resolution
digital scans of classical paintings (30 oil, 15 ink-wash, 5 fresco)
spanning Western and East Asian traditions. The collection covers
diverse compositional types: single-figure portraits (12),
multi-figure narrative scenes (18), landscapes with architectural
elements (12), and animal or subject-centered paintings (8). Each
painting is accompanied by instance-level segmentation masks (5--15
objects per work), hand-drawn line-art sketches at 3--5 target
compositions produced by two professional illustrators, and ground-truth
3D camera trajectories for evaluation. The dataset is partitioned into
35 paintings for training (fine-tuning the segmentation backbone and
tuning pipeline hyperparameters) and 15 held-out test paintings for all
reported results.

\paragraph{Implementation.}
All components are implemented in PyTorch. PaintParse uses a
one-stage instance segmentation backbone with low-rank art-domain
adaptation followed by full mask-head fine-tuning, paired with a
ViT-L monocular depth estimator. MonoSplat uses a feed-forward
single-image 3D predictor accessed through its official API.
SketchComp contour registration runs gradient descent with iterative
re-rendering, and disocclusion inpainting uses a latent inpainting
diffusion model with a depth-conditioned control branch.
ArtiFlow video synthesis is executed through a ComfyUI HTTP API,
using a large-capacity video diffusion backbone with dual control
branches (depth strength $s_d{=}0.85$, edge strength $s_e{=}0.6$) and
first-frame latent anchoring. The keyframe count is fixed to
$T=25$ (i.e., $8N{+}1$ with $N{=}3$).

\paragraph{Hardware.}
Experiments run on a single 48\,GB GPU. Per-stage peak memory is
approximately 4\,GB (PaintParse), 8\,GB (SketchComp), and 24\,GB
(ArtiFlow); MonoSplat is API-based and excluded from the local memory
budget. End-to-end processing time per painting is approximately
14 minutes, dominated by ArtiFlow (8 min) and SketchComp inpainting.

\paragraph{Baselines.}
We compare against six methods spanning the relevant paradigms.
\emph{Direct I2V} baselines include two state-of-the-art video diffusion
backbones applied to the full painting with default parameters.
\emph{Generative-prior clipart animation} represents video score
distillation from text-to-video priors, adapted for painting inputs.
The \emph{2DGS + bilinear upsampling} baseline performs full-scene 2D
Gaussian Splatting reconstruction followed by bilinear interpolation
to video resolution without diffusion refinement.
\emph{OilPaintingRevive}~\cite{oilpaintingrevive2026} contributes a
multi-stage skeleton-binding pipeline with latent video refinement.
\emph{LivePortrait}~\cite{guo2024liveportrait} is evaluated on the
portrait subset (12 paintings) only, applying face animation to the
dominant figure in each work.

\paragraph{Metrics.}
We report \textbf{CLIP-I} (CLIP image cosine similarity between
generated frames and the original painting) for style fidelity,
\textbf{LPIPS}~\cite{zhang2018lpips} for perceptual quality,
\textbf{FID-VID} (Fr\'echet Video Distance~\cite{balaji2019fvd}
computed with a video-trained I3D backbone) for temporal coherence,
\textbf{Disocclusion PSNR} for inpainted gap quality against a
manually-inpainted reference, and \textbf{Composition IoU} for the
average bounding-box overlap between detected objects in generated
frames and the target positions specified by the line-art sketches
(measuring compositional-intent alignment).

\subsection{Main Results (Q1)}

Table~\ref{tab:main_results} reports the primary quantitative
comparison. Canvas3D achieves the highest style fidelity (CLIP-I 0.86)
and the lowest perceptual distance (LPIPS 0.28), outperforming the
next-best method by 3.6\% and 9.7\% respectively. Its temporal
coherence (FID-VID 18.3) improves by 19.7\% over the next-best
painting-specific baseline and by 35.6\% over the strongest direct-I2V
baseline. Notably, Canvas3D is the only method with meaningful
Composition IoU (0.81), as no baseline supports object repositioning
per external line-art. The pure 2DGS baseline achieves the lowest
CLIP-I (0.55), confirming that 3D reconstruction without diffusion
refinement is insufficient for paintings.

\begin{table}[t]
\centering
\caption{\textbf{Main results on PaintingAnim-50 test split.}
Bold: best. ``--'' denotes metric inapplicable.}
\label{tab:main_results}
\resizebox{\linewidth}{!}{%
\begin{tabular}{@{}lccccc@{}}
\toprule
\textbf{Method} & \textbf{CLIP-I}$\uparrow$ & \textbf{LPIPS}$\downarrow$
    & \textbf{FID-VID}$\downarrow$ & \textbf{PSNR}$\uparrow$ &
    \textbf{IoU}$\uparrow$ \\
\midrule
Direct I2V (backbone A)        & 0.79 & 0.35 & 28.4 & --    & --    \\
Direct I2V (backbone B)        & 0.76 & 0.38 & 31.2 & --    & --    \\
Generative-prior clipart anim. & 0.72 & 0.41 & 34.7 & --    & --    \\
2DGS + bil.\ ups.              & 0.55 & 0.52 & 45.1 & --    & --    \\
OilPaintingRevive~\cite{oilpaintingrevive2026}
                               & 0.83 & 0.31 & 22.8 & --    & --    \\
LivePortrait$^*$~\cite{guo2024liveportrait}
                               & 0.74 & 0.39 & 30.5 & --    & --    \\
\midrule
\textbf{Canvas3D}              & \textbf{0.86} & \textbf{0.28} & \textbf{18.3}
    & \textbf{29.1} & \textbf{0.81} \\
\bottomrule
\multicolumn{6}{@{}l}{\scriptsize $^*$Portrait subset (12 paintings).}
\end{tabular}}
\end{table}

Figure~\ref{fig:qual_main} provides qualitative side-by-side
comparisons. Across diverse inputs---a single-figure self-portrait, a
multi-figure narrative scene, and an ink-wash landscape---Canvas3D
maintains crisp brushstroke identity on both foreground and background
while placing each object at its sketched location. Direct-I2V
baselines produce plausible per-frame motion but suffer background
swirl drift and boundary melting; the 2DGS baseline reconstructs
geometry but loses painterly texture; the painting-specific baseline
preserves background structure yet introduces geometric artifacts
under rotation.

\begin{figure}[t]
    \centering
    \fbox{\parbox[c][4.2cm][c]{0.97\linewidth}{\centering
        \textit{[Figure placeholder: Qualitative comparison.}\\[2pt]
        \textit{Three rows (self-portrait / multi-figure scene /
        ink-wash landscape) $\times$ columns: Input painting,
        Direct-I2V (backbone A), 2DGS baseline,
        OilPaintingRevive~\cite{oilpaintingrevive2026}, Canvas3D
        (ours). Arrows indicate object repositioning per sketch.
        Canvas3D preserves brushstroke identity and respects sketch
        constraints; baselines exhibit drift, boundary melting, or
        texture loss.]}}}
    \caption{\textbf{Qualitative comparison on three representative
    paintings.} Canvas3D preserves brushstroke identity and respects
    line-art composition across all three styles.}
    \label{fig:qual_main}
\end{figure}

\subsection{Ablation Studies (Q2)}

We isolate each module's contribution (Table~\ref{tab:ablations}).

\paragraph{PaintParse: art-domain adaptation.}
Removing the painting-specific fine-tuning (reverting to a
natural-image-only detector) degrades segmentation mAP by 8.3\%. This
cascades downstream: CLIP-I drops by 5.8\% and Composition IoU falls
by 9.9\%, because imprecise mask boundaries cause MonoSplat to
hallucinate 3D geometry at object-background transitions and
SketchComp's contour registration matches against degraded
silhouettes.

\paragraph{MonoSplat: feed-forward 3D vs.\ 2DGS back-projection.}
Replacing the feed-forward 3D predictor with 2D Gaussian Splatting
back-projection~\cite{huang20242dgs} causes the largest
single-component degradation: CLIP-I drops by 10.5\% and FID-VID
increases by 52.5\%. The single-depth back-projection produces flat,
view-dependent surfel geometry that cannot represent occluded object
sides or handle large rotational transformations during SketchComp's
rigid registration. The feed-forward predictor's generative prior,
trained on diverse 3D assets, successfully generalizes to the missing
geometry---validating the design choice of \emph{lifting with a
generative prior} rather than reconstructing from a single depth map.

\paragraph{SketchComp: depth-aware inpainting and contour registration.}
Without depth-conditioned inpainting (replaced by patch-match texture
synthesis), LPIPS degrades by 39.3\% and Disocclusion PSNR falls to
21.4\,dB. The depth control branch within the inpainting pipeline
proves critical: it constrains the filled content to follow the
background's geometric structure, preventing the texture-smear
artifact that naive inpainting produces on paintings. Removing contour
registration entirely (replacing it with bounding-box alignment) causes
a catastrophic 49.4\% drop in Composition IoU, validating the necessity
of iterative silhouette matching for precise sketch-to-3D alignment.

\paragraph{ArtiFlow: dual control and latent anchoring.}
Removing the edge control branch increases FID-VID by 52.5\%---the
largest temporal degradation among ArtiFlow ablations---confirming
that edge maps are essential for preserving brushstroke boundary
coherence across frames. Removing the depth control branch produces a
78.1\% FID-VID increase and a 14.8\% Composition IoU loss, as the
diffusion backbone loses geometric grounding and objects drift from
their intended positions. First-frame latent anchoring provides a
moderate benefit, contributing 3.5\% to CLIP-I and 31.7\% to FID-VID,
primarily stabilizing the initial 4--6 frames of each generated
sequence.

\begin{table}[t]
\centering
\caption{\textbf{Ablation results.} Each row removes one component.}
\label{tab:ablations}
\resizebox{\linewidth}{!}{%
\begin{tabular}{@{}lcccc@{}}
\toprule
\textbf{Variant} & \textbf{CLIP-I} & \textbf{LPIPS}
    & \textbf{FID-VID} & \textbf{IoU} \\
\midrule
Full Canvas3D                  & 0.86 & 0.28 & 18.3 & 0.81 \\
\midrule
-- Art adapt.\ (PaintParse)     & 0.81 & 0.34 & 23.7 & 0.73 \\
-- Feed-forward 3D (MonoSplat)  & 0.77 & 0.37 & 27.9 & 0.68 \\
-- Depth inpaint (SketchComp)   & 0.79 & 0.39 & 26.4 & 0.71 \\
-- Edge control (ArtiFlow)      & 0.82 & 0.31 & 27.9 & 0.76 \\
-- Depth control (ArtiFlow)     & 0.80 & 0.35 & 32.6 & 0.69 \\
-- Latent anchor (ArtiFlow)     & 0.83 & 0.30 & 24.1 & 0.78 \\
-- Contour reg.\ (SketchComp)   & 0.81 & 0.32 & 22.5 & 0.41 \\
\bottomrule
\end{tabular}}
\end{table}

\subsection{User Study (Q3)}

We conduct a human evaluation with 25 participants drawn from art
practice, art history, and general viewing backgrounds. Each
participant views 20 triplets consisting of a 5-second Canvas3D output
alongside outputs from the strongest direct-I2V baseline and
OilPaintingRevive~\cite{oilpaintingrevive2026} for the same painting,
rating each video on three 5-point Likert scales:
\textbf{artistic authenticity} (``How well does the video preserve the
painting's artistic character?''), \textbf{motion naturalness} (``How
natural and physically plausible is the object motion?''), and
\textbf{compositional intent} (``How well does the result match the
intended composition shown in the reference sketch?'').

Results (Table~\ref{tab:user_study}) show Canvas3D leading on all
three dimensions. The artistic-authenticity advantage over
OilPaintingRevive (mean 4.31 vs.\ 3.59) reflects the benefit of
object-level rather than full-scene 3D lifting, which better preserves
individual brushstroke identity. The compositional-intent gap is the
largest (4.48 vs.\ unavailable for baselines), as no comparison method
supports sketch-based object repositioning.

\begin{table}[t]
\centering
\caption{\textbf{User study} (Likert 1--5, mean $\pm$ std). N/A: cannot
accept line-art composition input.}
\label{tab:user_study}
\resizebox{\linewidth}{!}{%
\begin{tabular}{@{}lccc@{}}
\toprule
\textbf{Method} & \textbf{Art.\ Auth.} & \textbf{Motion Nat.}
    & \textbf{Comp.\ Intent} \\
\midrule
Direct I2V (backbone A) & 3.12 $\pm$ 1.05 & 3.48 $\pm$ 0.98 & N/A \\
OilPaintingRevive~\cite{oilpaintingrevive2026}
                        & 3.59 $\pm$ 0.87 & 3.72 $\pm$ 0.83 & N/A \\
\textbf{Canvas3D}       & \textbf{4.31}$\pm$\textbf{0.72} &
                          \textbf{4.05}$\pm$\textbf{0.79} &
                          \textbf{4.48}$\pm$\textbf{0.65} \\
\bottomrule
\end{tabular}}
\end{table}

\subsection{Compute Analysis (Q4)}

Table~\ref{tab:vram} details the per-stage memory budget and runtime.
ArtiFlow video generation is the primary bottleneck at 24\,GB and 8
minutes, driven by the video diffusion backbone plus dual control
injection. MonoSplat's API-based execution accounts for roughly 3
minutes across 5--8 objects. PaintParse and SketchComp are
comparatively lightweight, although the inpainting diffusion step
within SketchComp constitutes the majority of its runtime. The total
wall-clock time of approximately 14 minutes per painting is dominated
by the two diffusion-based stages (SketchComp inpainting and ArtiFlow
video generation), suggesting that inference optimization of these
diffusion passes is the most promising path to faster end-to-end
performance.

\begin{table}[t]
\centering
\caption{\textbf{Per-stage compute budget.} MonoSplat is API-based,
excluded from local memory.}
\label{tab:vram}
\resizebox{\linewidth}{!}{%
\begin{tabular}{@{}lccc@{}}
\toprule
\textbf{Stage} & \textbf{Memory} & \textbf{Runtime} & \textbf{Operation} \\
\midrule
PaintParse  & 4\,GB  & 0.5\,min  & Seg.\ + depth estimation \\
MonoSplat   & --     & 3.0\,min  & Per-object 3D lifting (API) \\
SketchComp  & 8\,GB  & 2.5\,min  & Registration + inpainting \\
ArtiFlow    & 24\,GB & 8.0\,min  & Video diffusion + dual control \\
\midrule
\textbf{Total} & \textbf{24\,GB} & \textbf{14\,min} & \\
\bottomrule
\end{tabular}}
\end{table}

\subsection{Viewpoint-Controllable Exploration (Q5)}

We evaluate ViewCtrl by generating five viewpoint trajectories per
test painting---one per mode (manual, orbit, dolly, crane,
auto-orbit)---and measuring two quantities:
\textbf{Novel-View IoU}, the average per-frame IoU between the
ViewCtrl-rendered foreground masks and the masks re-detected in the
SpatialFlow-generated frames (measuring whether the generated video
respects the viewpoint-specified visibility), and
\textbf{View-Consistency SSIM}, the average SSIM between consecutive
generated frames under viewpoint change (measuring temporal coherence
under novel-view exploration).

Table~\ref{tab:viewpoint} reports the results. Canvas3D-v2 maintains
a Novel-View IoU of 0.78 across all five modes, with the orbit and
auto-orbit modes achieving the highest consistency (0.82 and 0.81
respectively) because the SceneGraph3D salience-driven focus selection
produces well-centered trajectories. The crane mode is the most
challenging (0.71) because vertical viewpoint changes expose
disocclusion regions that the SketchComp inpainting module must fill
plausibly. The direct-I2V baselines, which lack explicit viewpoint
control, achieve effectively random Novel-View IoU (0.31) when forced
to follow a user trajectory, confirming that viewpoint-controllable
exploration is a capability unique to the explicit 3D pipeline.

\begin{table}[t]
\centering
\caption{\textbf{Viewpoint-controllable exploration (Q5).} Novel-View
IoU and View-Consistency SSIM per trajectory mode. Direct-I2V baselines
lack viewpoint control.}
\label{tab:viewpoint}
\resizebox{\linewidth}{!}{%
\begin{tabular}{@{}lcc@{}}
\toprule
\textbf{Mode} & \textbf{Novel-View IoU}$\uparrow$ &
    \textbf{View-Consist.\ SSIM}$\uparrow$ \\
\midrule
Manual (user poses)             & 0.76 & 0.84 \\
Orbit (focus node)              & 0.82 & 0.87 \\
Dolly (main axis)               & 0.78 & 0.85 \\
Crane (vertical boom)           & 0.71 & 0.81 \\
Auto-orbit (salience focus)     & 0.81 & 0.86 \\
\midrule
Direct I2V (no viewpoint ctrl)  & 0.31 & 0.62 \\
\bottomrule
\end{tabular}}
\end{table}

\subsection{Spatial-Relation Accuracy (Q6)}

We evaluate SpatialFlow by constructing 30 spatial-relation edits
across the 15 test paintings (e.g.\ ``move person left-of building,
behind tree''), generating videos, and measuring
\textbf{Relation Accuracy}: the fraction of generated frames in which
the edited relation holds, as detected by re-running PaintParse on
each frame and checking the relation against the edited graph.
We also report the \textbf{Re-prompt Rate}: the fraction of edits
that triggered the SpatialFlow verification loop at least once.

Table~\ref{tab:relation} shows that SpatialFlow achieves a Relation
Accuracy of 0.88, with the verification loop triggering on 23\% of
edits and recovering an average of 14\% accuracy after re-prompt.
Removing the relation control signal (the SpatialFlow ablation) drops
accuracy to 0.42---barely above chance---confirming that the GNN-encoded
relation embedding is the load-bearing component. Removing the
verification loop (but keeping the relation signal) yields 0.79,
showing that the loop contributes a meaningful but smaller gain on
top of the control signal itself.

\begin{table}[t]
\centering
\caption{\textbf{Spatial-relation accuracy (Q6).} Relation Accuracy
and Re-prompt Rate across configurations.}
\label{tab:relation}
\resizebox{\linewidth}{!}{%
\begin{tabular}{@{}lcc@{}}
\toprule
\textbf{Configuration} & \textbf{Relation Acc.}$\uparrow$ &
    \textbf{Re-prompt Rate}$\downarrow$ \\
\midrule
Full SpatialFlow (relation + verify)         & \textbf{0.88} & 0.23 \\
-- Verification loop (relation only)         & 0.79 & 0.00 \\
-- Relation signal (verify only)             & 0.48 & 0.61 \\
-- Both (depth + edge only, = ArtiFlow)      & 0.42 & 0.00 \\
\bottomrule
\end{tabular}}
\end{table}

\subsection{Ablation: v2 Modules (Q2 continued)}

We additionally ablate the three v2 modules
(Table~\ref{tab:ablations_v2}). Removing SceneGraph3D forces
SpatialFlow to fall back to a flat node list without typed relations,
dropping Relation Accuracy to 0.52 and increasing FID-VID by 18.6\%.
Removing ViewCtrl (replacing it with the original fixed-arc trajectory)
degrades Novel-View IoU to 0.39 and FID-VID by 12.0\%, as the
fixed-arc cannot respect user viewpoint specifications. Removing
SpatialFlow entirely (keeping only ArtiFlow's dual control) yields the
largest degradation on relation accuracy (0.42) and a 9.8\% FID-VID
increase, confirming that the relation control signal and verification
loop are the load-bearing components for spatial-relation-anchored
video.

\begin{table}[t]
\centering
\caption{\textbf{v2 module ablations.} Each row removes one v2 module.}
\label{tab:ablations_v2}
\resizebox{\linewidth}{!}{%
\begin{tabular}{@{}lcccc@{}}
\toprule
\textbf{Variant} & \textbf{CLIP-I} & \textbf{FID-VID} &
    \textbf{NV-IoU} & \textbf{Rel-Acc} \\
\midrule
Full Canvas3D-v2                 & 0.86 & 18.3 & 0.78 & 0.88 \\
\midrule
-- SceneGraph3D (flat nodes)     & 0.84 & 21.7 & 0.71 & 0.52 \\
-- ViewCtrl (fixed arc)          & 0.85 & 20.5 & 0.39 & 0.74 \\
-- SpatialFlow (depth+edge only) & 0.83 & 20.1 & 0.69 & 0.42 \\
\bottomrule
\end{tabular}}
\end{table}

\subsection{Additional Qualitative Results}

Figure~\ref{fig:qual_pipeline} visualizes the intermediate
representations produced by each stage of Canvas3D on a representative
input, illustrating how the handshake representations---masks with
depth-ordered layers, the 3D object library, the composited scene with
disocclusion filling, and the multi-view keyframe stream---progressively
convert a static painting into an explorable, artistically controllable
animation.

\begin{figure}[t]
    \centering
    \fbox{\parbox[c][3.5cm][c]{0.97\linewidth}{\centering
        \textit{[Figure placeholder: Pipeline visualization.}\\[2pt]
        \textit{Left to right: Input painting $\to$ PaintParse masks
        with depth-ordered layer coloring $\to$ MonoSplat 3D library
        (rendered from a novel view) $\to$ SketchComp composited scene
        with disocclusion gaps filled $\to$ ArtiFlow multi-view
        keyframes $\to$ final animated frames.]}\\[2pt]
        \textit{Each transition shows the handshake representation
        passed between stages.]}}}
    \caption{\textbf{Pipeline visualization.} Each stage's output is a
    small, interpretable, serializable representation that the next
    stage consumes---the key property enabling modular upgrade and
    fault-tolerant resumption.}
    \label{fig:qual_pipeline}
\end{figure}

Figure~\ref{fig:qual_novelview} shows novel-view exploration of the
composited 3D scene, demonstrating that the framework supports
viewpoint changes beyond the original painting's perspective while
preserving brushstroke identity---a capability no pure 2D I2V method
can offer.

\begin{figure}[t]
    \centering
    \fbox{\parbox[c][3.2cm][c]{0.97\linewidth}{\centering
        \textit{[Figure placeholder: Novel-view exploration.}\\[2pt]
        \textit{Top row: original viewpoint and four rendered novel
        viewpoints of the composited 3D scene.}\\
        \textit{Bottom row: corresponding frames from the generated
        video, showing temporally consistent novel-view motion.]}}}
    \caption{\textbf{Novel-view exploration.} Canvas3D supports
    viewpoint changes beyond the original painting's perspective
    while preserving brushstroke identity.}
    \label{fig:qual_novelview}
\end{figure}

Figure~\ref{fig:qual_inpainting} compares depth-conditioned inpainting
against naive (patch-match) inpainting on the disocclusion gaps opened
by object repositioning, illustrating the texture-smear artifact that
the depth control branch is designed to prevent.

\begin{figure}[t]
    \centering
    \fbox{\parbox[c][3.0cm][c]{0.97\linewidth}{\centering
        \textit{[Figure placeholder: Disocclusion inpainting
        comparison.}\\[2pt]
        \textit{Columns: Original painting $\to$ Object repositioned
        (gap visible) $\to$ Naive patch-match inpainting (texture
        smear) $\to$ Depth-conditioned inpainting (geometry-aware
        fill) $\to$ Ground-truth manual inpainting.]}}}
    \caption{\textbf{Depth-conditioned inpainting.} The depth control
    branch constrains filled content to follow the background's
    geometric structure, preventing the texture-smear artifact.}
    \label{fig:qual_inpainting}
\end{figure}
